%% toc: P15 164 164.2
% Source-specific presentation: J. R. Milton, Induction before Hume (HHL 10, ch. 1, pp. 1--41).
\chapter{Induction before Hume}
\label{sec:p15-w02b-p15-026-milton}

This chapter reconstructs J.~R. Milton's history of induction from
the ancient world to Leibniz, the opening chapter of the
\emph{Inductive Logic} volume of the \emph{Handbook of the History of
Logic}. Its purpose is to let a reader locate any pre-Humean remark
about induction: which word the writer used, what the argument was
supposed to lead to (a concept, a universal proposition or another
particular), what certainty it was claimed to give, and how Milton
reads it against its predecessors. The input is a body of scattered
remarks by writers who were mostly occupied with other matters. The
conclusions are Milton's historical theses. Induction was named and
first analysed by Aristotle but never given a treatment comparable to
his syllogistic. The fallibility of incomplete enumeration was a
commonplace from Sextus and Alexander of Aphrodisias onwards. Bacon is
the one writer for whom induction was central. No one before the
nineteenth century saw a single general problem of induction
\cite[pp.~1--2, 8, 10, 36--37]{Milton2011InductionBeforeHume}.

Page locators refer to Milton's chapter, pp.~1--41, read in the
camera-ready text of the volume (copyright footer 2009, Elsevier BV).
Milton's own references to primary texts (Bekker numbers, book and
chapter numbers, columns of the 1608 Zabarella edition, pages of the
translations he quotes) are reported as he gives them, so that a reader
can follow him to his sources; they are his locators, not new ones.
Hume's own argument is treated in
\cref{sec:p15-hume-and-the-problem-of-induction}, and Mill's methods,
which Milton invokes in his first example, in \cref{sec:p15-wm-methods}.

\begin{hbconvention}{def:p15-w02b-p15-026-claim-types}{Kinds of claim in this chapter}
Four kinds of claim are kept apart.
\begin{enumerate}
\item \emph{Textual-historical}: what a writer wrote, when, and in which
  work, as Milton reports it; and Milton's interpretation of it, marked as
  his.
\item \emph{Philosophical}: theses the historical writers argue for, for
  example that demonstrative induction yields certainty, or that methodic
  experience does and induction does not.
\item \emph{Mathematical}: statements about the explicitly defined modern
  reconstructions in this chapter (the lemmas below). They are proved in the
  annex. They describe what follows in the reconstruction under its stated
  hypotheses; they are not claims about what the historical writer proved or
  meant.
\item \emph{Empirical}: observations the historical writers report (the
  purgative action of scammony and rhubarb, salt seas and fresh rivers,
  syringes, suction pumps and narrow tubes). Milton reports them as examples
  in historical texts. The chapter gives no population, design, measures or
  uncertainty for any of them, and nothing here is a new empirical finding.
\end{enumerate}
\end{hbconvention}

\section{The words and the extent of the history}
\label{sec:p15-w02b-p15-026-terms}

\begin{hbdefinition}{def:p15-w02b-p15-026-terms}{Epagoge, inductio, istiqra' and semeiosis}
The following terms, as Milton uses and glosses them, organize the history.
\begin{enumerate}
\item \emph{Epagoge} is Aristotle's name for inductive reasoning, formed
  from the verb \emph{epagein}, rendered by context as to bring or lead
  in, or on. The noun had unrelated senses as well: an incantation in
  Plato (\emph{Republic} 364c) and the taking in of food in Aristotle
  (\emph{De Respiratione} 483a9)
  \cite[p.~2]{Milton2011InductionBeforeHume}.
\item \emph{Inductio} is the Latin word used for \emph{epagoge}. Who
  coined the equivalence is unknown; in the surviving Latin corpus the
  word first appears in this sense in Cicero's early \emph{De Inventione}
  \cite[p.~11]{Milton2011InductionBeforeHume}. Modern ``induction'' and
  its cognates descend from it and also carry irrelevant senses
  (electromagnetic induction, the induction of a cleric into a benefice)
  \cite[p.~1]{Milton2011InductionBeforeHume}.
\item \emph{Istiqra'} is the Arabic rendering of \emph{epagoge}, a word
  whose root sense is investigation or examination
  \cite[p.~14]{Milton2011InductionBeforeHume}.
\item \emph{Semeiosis}, inference from signs, was the Hellenistic centre
  of discussion of non-deductive inference. It argues from the apparent to
  the non-apparent, where the latter is either temporarily non-apparent
  (something at a distance) or permanently and intrinsically so (invisible
  bodies such as atoms) \cite[p.~1]{Milton2011InductionBeforeHume}.
\end{enumerate}
\end{hbdefinition}

Milton states two methodological decisions at the outset. First, a history
confined to the direct ancestors of present usage would be hard to write
and unsatisfying if finished, so he casts the net wider than the modern
sense and includes uses that diverge from it. Second, the vocabulary is not
a sufficient guide in the other direction either: a writer can discuss
inductive inference without the word. Hume is the main case. He is widely
read, at least in the twentieth century, as an inductive sceptic although
he seldom used the word and never in the passages where that scepticism
is located; Milton's note~1 cites Stove (1973), Winkler (1999), Howson
(2000) and Okasha (2001) for that reading. He leaves its correctness
aside but notes
that it cannot be settled by examining Hume's vocabulary
\cite[p.~1]{Milton2011InductionBeforeHume}.

\begin{hbwarning}{def:p15-w02b-p15-026-signs}{Sign inference and the word ``induction''}
Whether the vocabulary of induction should be applied to Hellenistic sign
inference is disputed among modern historians. Milton records two opposed
practices: Asmis (1984) uses the terminology freely, while Barnes (1988)
rejects it altogether. Milton does not adjudicate. He declines to discuss
sign inference in detail and refers to Burnyeat (1982), Sedley (1982) and
Allen (2001) for modern accounts
\cite[p.~1]{Milton2011InductionBeforeHume}. The present chapter follows
him: nothing below reconstructs Hellenistic sign inference, and a reader who
needs it must go to those studies. Calling a Hellenistic sign inference
``inductive'' already takes a side in this dispute.
\end{hbwarning}

\section{The ancient world}
\label{sec:p15-w02b-p15-026-ancient}

\subsection{Before the logicians}

Generalization long preceded its analysis, and practice could exceed
induction by simple enumeration. Milton's first witness is a narrative,
not a theory.

\begin{hbexample}{def:p15-w02b-p15-026-gideon}{Gideon's fleece (Judges vi.~36--40)}
Gideon asks for a sign that he will save Israel. He sets a fleece on the
threshing floor and stipulates that dew should fall on the fleece alone
while the surrounding ground stays dry; in the morning he wrings a bowl of
water from the fleece. He then asks to try once more with the conditions
reversed, the fleece dry and the ground wet, and that night it happens
so. Milton quotes the Authorized Version
\cite[pp.~1--2]{Milton2011InductionBeforeHume}.

Milton's comment is that neither the writer nor the original readers knew
Mill's Methods of Agreement or Difference, yet few would have found
Gideon's procedure hard to follow. He adds Locke's sardonic remark that God
did not leave it to Aristotle to make men rational (\emph{Essay}
IV.~xvii.~4) \cite[p.~2]{Milton2011InductionBeforeHume}.

What makes the procedure more than enumeration is the second trial.
The following gloss is this chapter's, not Milton's. The first night
alone is compatible with the hypothesis that a fleece simply gathers
more dew than bare ground. The reversed trial varies the condition
under which the pattern would follow from that ordinary hypothesis and
obtains the opposite outcome. That is the contrastive structure which
Mill later called difference (\cref{def:p15-wm-difference}). The
narrative supplies no sample, controls or measures beyond the two
nights.
\end{hbexample}

Aristotle is nevertheless the first to name inductive reasoning and to
give an account of it, brief and imperfect as it is
\cite[p.~2]{Milton2011InductionBeforeHume}.

\subsection{Socrates and Plato}

Aristotle did not inherit the word \emph{epagoge} for these arguments, but
he singled out Socrates for using what he called \emph{epaktikous logous},
inductive arguments (\emph{Metaphysics} 1078b28). Milton takes Plato to be
Aristotle's main source on Socrates
\cite[p.~2]{Milton2011InductionBeforeHume}.

\begin{hbexample}{def:p15-w02b-p15-026-protagoras}{Opposites in the \emph{Protagoras} (332C)}
Socrates secures from Protagoras, one case at a time, that the beautiful
has only the ugly as its opposite, the good only the evil, and the acute in
sound only the grave. He then asks whether every opposite has exactly one
opposite, and Protagoras agrees \cite[pp.~2--3]{Milton2011InductionBeforeHume}.

Milton draws the boundary that matters for the rest of the history. Here,
and in \emph{Charmides} 159--160 and \emph{Ion} 540, the conclusion is
philosophical: a sufficiently clear-sighted person could grasp it directly.
Plato's concern was with truths of this kind, not with empirical
generalizations about sensible particulars such as white swans
\cite[p.~3]{Milton2011InductionBeforeHume}. The cases elicit assent; they
are not a sample from which a natural regularity is extrapolated.
\end{hbexample}

\subsection{Aristotle}
\label{sec:p15-w02b-p15-026-aristotle}

Milton writes ``epagoge'' rather than ``induction'' because even the
translation is disputed. He quotes Upton (1981) on the common belief that
Aristotle's concept is incomplete and unsystematic compared with his theory
of deduction, and concludes that there is no consensus on what Aristotle
attempted or how far he succeeded. His note lists a range of views:
Kosman, Hamlyn, Engberg-Pedersen, Upton, Caujolle-Zaslavsky, McKirahan and
De Rijk \cite[p.~3]{Milton2011InductionBeforeHume}.

When Aristotle applies \emph{epagoge} to his own arguments the usage is
Socratic: the conclusions are seldom empirical generalizations. The typical
case is \emph{Metaphysics} 1055a5--10, where induction is said to show that
contrariety is the greatest difference. Things differing in genus are too
far apart to be compared; for things differing in species, generation
proceeds between contraries as extremes, and the distance between extremes
is greatest. Similar uses occur at 1055b17 and 1058b9
\cite[p.~3]{Milton2011InductionBeforeHume}.

Aristotle discusses \emph{epagoge} in three passages, none long.

\begin{hbdefinition}{def:p15-w02b-p15-026-topics}{The \emph{Topics} characterization}
Dialectical arguments are of two kinds, \emph{syllogismos} (best rendered
``deduction'', a wider notion than the later syllogism) and
\emph{epagoge}. \emph{Epagoge} is the progress from particulars to
universals. In the example, the best pilot and the best charioteer are each
the skilled one, so in general, in any given sphere, the skilled person is
the best. Induction is the more convincing and clear,
more accessible to sense perception and shared by most people;
\emph{syllogismos} is the more compelling and the more effective against
opponents in argument \cite[p.~4]{Milton2011InductionBeforeHume}.
\end{hbdefinition}

The first clause became the standard medieval and Renaissance definition.
Milton qualifies its modern reading. It is natural to take it as saying
that induction is inference from particular to universal propositions,
but the Greek says only that induction is the passage (\emph{ephodos})
from individuals to universals (\emph{ta katholou}). Elsewhere, notably in
\emph{Posterior Analytics} B19, these universals appear to be or to include
universal concepts, and \emph{ephodos} should not be assumed to mean
inference in a technical sense (citing De Rijk)
\cite[p.~4]{Milton2011InductionBeforeHume}.

\begin{hbexample}{def:p15-w02b-p15-026-bileless}{The bileless animals (\emph{Prior Analytics} B23, 68b15--29)}
Aristotle's longest account calls the procedure ``the syllogism which
springs out of induction'' (\emph{ho ex epagoges syllogismos}). It
establishes a relation between one extreme and the middle term by means of
the other extreme: if $B$ is the middle between $A$ and $C$, one proves
through $C$ that $A$ belongs to $B$. The letters are assigned as follows:
\begin{itemize}
\item $A$: long-lived;
\item $B$: bileless;
\item $C$: the particular long-lived animals, such as man, horse and mule.
\end{itemize}
$A$ belongs to all $C$, and $B$ belongs to all $C$. If $C$ is convertible
with $B$, the middle not being wider in extension, then $A$ must belong to
$B$. Aristotle appeals to a result already proved: when two things belong
to the same thing and the extreme converts with one of them, the other
predicate belongs to the converted one. He adds that $C$ must be
apprehended as made up of all the particulars
\cite[p.~4]{Milton2011InductionBeforeHume}.

Milton records a textual problem. In the Oxford rendering, a clause after
``$A$ belongs to the whole of $C$'' says that whatever is bileless is
long-lived. Milton says the clause makes no sense in that position; it
states the very conclusion the passage is proving (this chapter's
observation). The 1973 Cooke--Tredennick translation suggests it may be
an interpolation to be excised; Ross notes that no manuscript supports
excising it (Milton's note~3) \cite[p.~4]{Milton2011InductionBeforeHume}.
\end{hbexample}

Milton's reading of the passage has four steps
\cite[p.~5]{Milton2011InductionBeforeHume}. (i)~Aristotle is applying his
method of conversion, devised for the syllogistic, to a case it does not
obviously fit; hence the talk of middle terms. (ii)~The crucial premise is
that $B$ belongs to all $C$, that every long-lived animal is bileless.
This could range over individual animals or over species. The examples
(man, horse, mule, not named individuals) suggest species. (iii)~If so, the
example is what was later called perfect, that is complete, induction: an
inference from a finite sample small enough for every case to be examined.
(iv)~The closing phrase seems to say this, but the Oxford translation adds
words. The Greek (\emph{he gar epagoge dia panton}) says only that
induction is ``through all''.

\begin{hblemma}{thm:p15-w02b-p15-026-conversion}{proved=annex:p15-w02b-p15-026-conversion}{The conversion step in an extensional reconstruction}
In the following modern reconstruction each term is identified with
its extension, a set of things. ``$X$ belongs to all $Y$'' is read as
$Y\subseteq X$. ``$C$ is convertible with $B$'' is read as the two
inclusions $C\subseteq B$ and $B\subseteq C$. Under this reading:
\begin{enumerate}
\item from $C\subseteq A$, $C\subseteq B$ and $B\subseteq C$ it follows that
  $B\subseteq A$, that is, $A$ belongs to all $B$;
\item from $C\subseteq A$ and $C\subseteq B$ alone it does not follow that
  $B\subseteq A$.
\end{enumerate}
\end{hblemma}

The lemma locates the inductive content. Part~(1) is a valid deduction, and
it is valid whatever the size of $C$. Part~(2) shows that the inference
depends on the converse $B\subseteq C$. The converse says that everything
bileless is among the long-lived animals listed in $C$, and that is what
the requirement that $C$ be ``made up of all the particulars'' must
supply. The reconstruction does not decide whether $C$ ranges over
individuals or species, and it does not adjudicate the textual problem
above. Its relation to the modern syllogistic calculi is treated in
\cref{sec:p02-syllogistic-inference}.

Another proposition said to be reached by induction occurs in Aristotle's
celebrated example of scientific demonstration, \emph{Posterior Analytics}
A13 (78a30--b4)
\cite[p.~5]{Milton2011InductionBeforeHume}:
\begin{center}
\begin{tabular}{rl}
(1) & The planets do not twinkle.\\
(2) & Whatever does not twinkle is near.\\
therefore (3) & The planets are near.
\end{tabular}
\end{center}
It counts as a demonstration, not merely a valid syllogism, because it
gives the cause: the planets do not twinkle because they are nearer than
the fixed stars. Aristotle says premise (2) was reached by induction or
through perception (78a34--5). Milton observes that the same must hold of
(1). For (1) the induction is unproblematic: Mercury does not twinkle,
Venus does not, and so on through the planets. For (2) it is not. A long
list of nearby non-twinkling objects is easy to assemble, but if (2) is the
conclusion of an induction, the enumeration is plainly incomplete and the
inference correspondingly fallible. Like the bileless animals, and unlike
the Socratic cases, (2) is a generalization about nature that could not be
grasped by intuition \cite[p.~5]{Milton2011InductionBeforeHume}.

Demonstrations proceed from premises that are, in the Oxford rendering
Milton quotes, ``true, primary, immediate, better known than, prior to, and
causative of their conclusion'' (71b20--2).
All such premises are universal, so the question arises how the
indemonstrable first premises are known. Aristotle postpones it to the
last chapter of the second book (B19). There, after stating the problem,
he appears to describe the acquisition of universal concepts, from
the perception of a number of individual men to the species, and from
there to the genus animal (100a3--b3). The single occurrence of \emph{epagoge} in the chapter
then says that first principles must be known by induction, since this is
also how perception implants universals (100b3--5)
\cite[pp.~5--6]{Milton2011InductionBeforeHume}.

\begin{hbwarning}{def:p15-w02b-p15-026-ross-barnes}{Two readings of \emph{Posterior Analytics} B19}
Milton sets out two incompatible readings by the main English
commentators \cite[p.~6]{Milton2011InductionBeforeHume}.
\begin{itemize}
\item \emph{Ross (1949).} The chapter treats concept formation and
  induction together, because forming general concepts and grasping
  universal propositions are inseparably interwoven.
\item \emph{Barnes (1975).} ``Induction'' is used here in a weak sense for
  any cognitive progress from less to more general, and 100b3--5 says only
  that concept acquisition proceeds in that direction; the passage is not
  about inference from singular to universal propositions at all.
\end{itemize}
Milton judges that the dispute cannot easily be resolved: the texts are
short and all parties know them thoroughly. He inclines to Ross. His reason
is a contrast with Locke. For Locke, abstraction yields an abstract general
idea, a particular entity able to function as a universal sign, and the
process has nothing to do with an inductive ascent between propositions,
about which Locke says almost nothing. For Aristotle what comes to rest in
the intellect is a real universal, a form freed from matter and so
de-individuated. That is why one psychological process can explain both
concept acquisition and knowledge of first principles. Aristotle gives
little more than a sketch; his late ancient and medieval followers worked
it out \cite[p.~6]{Milton2011InductionBeforeHume}.
\end{hbwarning}

Milton adds a separating boundary. Aristotle gives no hint that
\emph{epagoge} is only one of several routes to first principles. The
familiar empiricist division, with some universal truths supported a
posteriori by induction and others, such as Euclid's equality of right
angles, known a priori, is foreign to him. At \emph{Posterior Analytics}
81b2 universals can be viewed only through induction
\cite[p.~6]{Milton2011InductionBeforeHume}.

\begin{hbdefinition}{def:p15-w02b-p15-026-paradeigma}{Example (\emph{paradeigma})}
Much practical reasoning runs from particulars to particulars: whether the
sun will rise tomorrow, whether this loaf will nourish. Aristotle classes
such inferences as examples, \emph{paradeigmata}
\cite[pp.~6--7]{Milton2011InductionBeforeHume}.
\begin{itemize}
\item \emph{Prior Analytics} B24 (68b41--69a13): Thebes fighting Phocis
  and Athens fighting Thebes are both wars on neighbours. Since the first
  turned out badly for Thebes, the second will turn out badly for
  Athens.
\item The inference may seem to pass through the principle that war
  against neighbours is always bad (69a4, 6). It would then be an inductive
  ascent followed by a deductive descent. Aristotle nonetheless separates the
  two. Example goes from part to part, not from part to whole or whole to
  part (69a14--15), and induction proceeds from all the individual cases
  while example does not (69a16--19)
  \cite[p.~7]{Milton2011InductionBeforeHume}.
\item \emph{Rhetoric} 1356b1--10: the example is a rhetorical induction and
  the enthymeme a rhetorical syllogism. Everyone who persuades by proof uses
  one or the other. The exhaustive division of arguments into
  \emph{syllogismos} and \emph{epagoge} also occurs in both
  \emph{Analytics} (68b13--14, 71a5--6), together with the pairing of
  enthymeme and example with them (71a9--11)
  \cite[p.~7]{Milton2011InductionBeforeHume}.
\end{itemize}
\end{hbdefinition}

Milton thus records that the \emph{Analytics} and the \emph{Rhetoric}
pull in different directions on whether example is a kind of induction.
One very plausible interpretation, Burnyeat's (1994), takes enthymeme and
example to be \emph{syllogismos} and \emph{epagoge} occurring in a
rhetorical speech rather than a dialectical argument. They are then
neither subspecies nor different kinds of argument. Milton states the
price: \emph{epagoge} must then be broadened to cover most or all
non-deductive argument, since \emph{paradeigma} certainly runs from
particulars to particulars. He declines to analyse the enthymeme and
refers to Burnyeat for it \cite[p.~7]{Milton2011InductionBeforeHume}.

No surviving work of Aristotle gives a detailed, systematic account of
induction, and there is no evidence one was written. Milton offers a
conjectural explanation, flagged as one possibility among others: the
separation of form and content, central to the analysis of the syllogism,
was and remains harder to achieve for induction. Aristotle therefore left
nothing comparable to his syllogistic \cite[p.~8]{Milton2011InductionBeforeHume}.

\subsection{Hellenistic and later Greek accounts}

In the three centuries after Aristotle his technical works were studied
little outside a declining Peripatetic school, and his terms were replaced.
Inference from particular to universal propositions was raised only
occasionally, unlike inference from signs
\cite[p.~8]{Milton2011InductionBeforeHume}.

\paragraph{Alcinous.}
The Platonist handbook (\emph{Didaskalikos}) ascribed to Alcinous, often
identified with the second-century Middle Platonist Albinus, gives the
whole treatment in two sentences. Induction is any procedure passing from
like to like or from particular to general, and it is especially useful
for activating the natural concepts (\emph{Didaskalikos} 6.7). Dillon
suggests an allusion to the slave-boy passage of the \emph{Meno}. Milton
reads the passage as a blend of Platonism and Aristotelianism: particular
examples bring to consciousness concepts implanted by nature
\cite[p.~8]{Milton2011InductionBeforeHume}.

\begin{hbdefinition}{def:p15-w02b-p15-026-diogenes}{Diogenes Laertius's division (\emph{Lives} III.~53--55)}
In the life of Plato, \emph{epagoge} is an argument inferring from some
true premises a conclusion resembling them. It has two varieties
\cite[pp.~8--9]{Milton2011InductionBeforeHume}.
\begin{enumerate}
\item \emph{From opposites.} The example: if man is not an animal he is a
  stick or a stone; he is neither, being animate and self-moved; so he is
  an animal. It continues: ``But if he is an animal, and if a dog or an ox is
  also an animal, then man by being an animal will be a dog and an ox as
  well.'' Milton finds the first part clear
  (an ancient ``animal, vegetable or mineral'') and the second obscure. It
  bears little resemblance to any modern notion of induction.
\item \emph{From implication}, with two sub-varieties: the rhetorical one,
  from particulars to other particulars, which is Aristotle's
  \emph{paradeigma} under another name; and the dialectical one, from
  particulars to universals. The dialectical instance is the
  \emph{Phaedo}'s proof of immortality. It rests on the general
  proposition that opposites come from opposites, which is itself
  established from particulars such as sleep from waking and the greater
  from the less, and conversely.
\end{enumerate}
Milton's comment on the second instance: these are not empirical
generalizations \cite[p.~9]{Milton2011InductionBeforeHume}.
\end{hbdefinition}

Milton characterizes Diogenes (early third century?) as a compiler rather
than a thinker, whose work survives as the only extensive ancient
compendium of philosophers' lives \cite[p.~8]{Milton2011InductionBeforeHume}.

\begin{hbexample}{def:p15-w02b-p15-026-sextus}{Sextus Empiricus's dilemma (\emph{Outlines of Pyrrhonism} II.~204 and II.~195)}
Inductive scepticism is inconspicuous in Sextus but present. In II.~204
he dismisses induction briefly. Those who establish a universal from
particulars must survey either all of the particulars or some. If some,
the induction is insecure, since an omitted particular may contravene the
universal. If all, they attempt the impossible, since particulars are
infinite and indefinite. Either way induction is ``invalidated''; Milton
notes that the word means literally shaken or made to totter
\cite[p.~9]{Milton2011InductionBeforeHume}.

II.~195 gives the target: ``every man is an animal'' is thought to be
established from Socrates, Plato, Dion and each other instance. A single
counter-example spoils such a universal, as when most animals move the
lower jaw and the crocodile alone the upper, so that ``every animal moves
the lower jaw'' is false \cite[p.~10]{Milton2011InductionBeforeHume}.
\end{hbexample}

Milton compares this with the modern black swan and finds the differences
small. What made the swan falsification effective was the discovery of a
\emph{species} of black swans. Logically one negative instance can falsify
a universal proposition. In practice a single instance usually would not,
because what Lakatos called monster-barring stratagems come into play. The
crocodile example was probably a stock one passed from writer to writer: it
also occurs in Apuleius' \emph{Peri Hermeneias}
\cite[p.~10]{Milton2011InductionBeforeHume}.

\begin{hblemma}{thm:p15-w02b-p15-026-enumeration}{proved=annex:p15-w02b-p15-026-enumeration}{Complete, covering and partial enumeration}
Let $D$ be a set of individuals, let $S\subseteq D$ be the extension of a
subject term and $P\subseteq D$ that of a predicate, and let $O\subseteq D$
be the set of examined individuals. Assume the evidence is
$S\cap O\subseteq P$: every examined individual falling under the subject
has the predicate.
\begin{enumerate}
\item \emph{Complete enumeration.} If $S\subseteq O$, then $S\subseteq P$.
\item \emph{Covering classes.} If $S=K_1\cup\dots\cup K_m$ for some
  $m\geq 1$ and $K_i\subseteq P$ for each $i$, then $S\subseteq P$. If an
  individual $x$ belongs to none of $K_1,\dots,K_m$, then $x\notin S$.
\item \emph{Partial enumeration.} If $S\not\subseteq O$, the evidence does
  not settle the universal. There are predicate extensions $P_1,P_2$ with
  $P_1\cap O=P_2\cap O=P\cap O$ such that $S\subseteq P_1$ and
  $S\not\subseteq P_2$.
\end{enumerate}
\end{hblemma}

The lemma is a modern reconstruction of the logical skeleton that recurs
from here to Leibniz, not a claim about any writer's own formulation.
Part~(3) is the first horn of Sextus's dilemma. Part~(1) is the second horn
\emph{with} the exhaustiveness premise $S\subseteq O$, which Sextus denies
can be secured when the particulars are infinite and indefinite. Part~(2)
is the form used by Aristotle's species reading, Gassendi and Descartes
below. The lemma says nothing about how probable the universal is after
partial enumeration. Milton records no numerical estimate of probability
in any writer he discusses, and of Wallis he says that no attempt was made
to estimate degrees of probability
\cite[p.~33]{Milton2011InductionBeforeHume}.

\paragraph{Alexander of Aphrodisias.}
That inductive conclusions are not conclusively established is not
peculiar to sceptics. Commenting on \emph{Topics} 105a10ff, the
second-century Aristotelian Alexander says induction has persuasiveness
but not necessity: the universal does not follow by necessity from the
conceded particulars, because the particular cases cannot be gone through
\cite[p.~10]{Milton2011InductionBeforeHume}.

Milton draws a historical conclusion from this. It is wholly mistaken to
suppose that Hume first noticed that inductive arguments are not
deductively valid, or that any generalization over an infinite or
unsurveyably large field is open to counter-example. The supposition would
be unfair to Hume, whose project was more radical and far less banal, and
to his predecessors, who took this fallibility for granted
\cite[p.~10]{Milton2011InductionBeforeHume}.

\subsection{Roman philosophy}

\paragraph{Cicero and the rhetorical tradition.}
The Romans had little interest in logic as a technical discipline,
though Milton notes, citing Barnes (1997), that the aversion was not
universal. Rhetoric, by contrast, was central to their education
\cite[p.~11]{Milton2011InductionBeforeHume}. In \emph{De Inventione}
induction is a technique of argument: the speaker first obtains agreement
to undisputed propositions and then leads the opponent to assent to
others resembling them.

\begin{hbexample}{def:p15-w02b-p15-026-aspasia}{Aspasia's questions (\emph{De Inventione} I.~55)}
Aspasia asks the wife of a certain Xenophon (not the historian) whether
she would prefer her neighbour's gold ornament, if it were better than her
own, and then the neighbour's finer clothes. She answers ``hers'' each
time. Asked whether she would prefer the neighbour's husband if he were
better than her own, she blushes \cite[p.~11]{Milton2011InductionBeforeHume}.

Milton classifies this as Aristotle's \emph{paradeigma}, not inductive
generalization. In the later \emph{Topics} (42) Cicero gives another case:
honesty is demanded of guardians, partners, bailees and trustees, and so
it is demanded of agents too. He calls it by the Greek name \emph{epagoge}
and treats it as one kind of argument from similarity, but it too
runs from particulars to a particular \cite[p.~11]{Milton2011InductionBeforeHume}.
\end{hbexample}

Cicero's concern was with argument in public speaking and the law courts,
not with the generalizations of a natural philosopher. Quintilian
(\emph{Institutio Oratoria} V.~x.~73, xi.~2) and later rhetoricians
followed him and treated induction cursorily, as less worth analysis than
signs or examples. Milton's assessment is that absorbing induction into
rhetoric meant the problem of how universal propositions can be proved,
or at least made probable, from particular cases was never properly
raised, let alone answered. On the meaning of \emph{probabilis} and
related terms in Cicero and other ancient writers his note~7 refers to
Glucker (1995), and on their later history to Hacking (1975), Cohen
(1980) and Franklin (2001) \cite[pp.~11--12]{Milton2011InductionBeforeHume}.

\paragraph{Boethius.}
Aristotle's works began to be translated into Latin only after the fall of
the western Empire. Of Boethius' planned translation of the whole corpus,
only the \emph{Categories} and \emph{De Interpretatione} survive, but his
own logical writings informed early medieval readers about the rest.
\emph{De Topicis Differentiis} defines induction in Aristotle's, not
Cicero's, terms as progression from particulars to universals. Its example
is adapted from the \emph{Topics}: choose a pilot, and likewise a
charioteer, for skill rather than by lot, and so in general choose by
skill whoever is to govern something. Milton judges that Boethius added
nothing to earlier analyses; his importance is that he supplied
information otherwise unavailable to his early medieval readers
\cite[p.~12]{Milton2011InductionBeforeHume}.

\paragraph{Milton's summary for antiquity.}
No sustained ancient discussion of induction survives. Much is lost, but
there is no sign that a systematic account has been lost. The lists of
works in Diogenes Laertius contain no treatise titled \emph{Peri Epagoges}
or similar \cite[p.~12]{Milton2011InductionBeforeHume}. Milton concludes
that no ancient writer treated induction as a central problem of
philosophy and gives three reasons \cite[pp.~12--13]{Milton2011InductionBeforeHume}.
\begin{enumerate}
\item Late ancient philosophy drifted away from systematic empirical
  inquiry of the kind Aristotle and his immediate successors practised.
  Plotinus uses \emph{epagoge} only twice: for an argument that nothing is
  contrary to substance, and for one that whatever is destroyed is
  composite (\emph{Enneads} I.~8.~6; II.~4.~6). Understanding gained by
  empirical generalization seemed too slight to merit analysis of the
  arguments that led to it.
\item Systematic investigation of nature was intermittent and local. Its
  most successful discipline, mathematical astronomy, does not generalize
  from ``this $A$ is $B$'' to ``every $A$ is $B$''. Ptolemy gives his models
  of the individual planets separately and, in the \emph{Almagest} at
  least, does not integrate them into one system.
\item The centrality of rhetoric meant inductive arguments were judged for
  persuasiveness rather than logical merit. They were lost among
  unformalized arguments assessed for plausibility in a speech, not for
  validity or any inductive analogue of validity.
\end{enumerate}

\section{The Middle Ages}
\label{sec:p15-w02b-p15-026-medieval}

\subsection{Arabic accounts}

From the late eighth century much of the philosophy fashionable in late
antiquity, including most of Aristotle and many commentators, was
translated into Arabic, and the accounts of induction in the
\emph{Analytics} and the \emph{Topics} became the starting point. Little
Greek technical vocabulary was transliterated (\emph{falsafah} is an
exception). \emph{Epagoge} became \emph{istiqra'}, with the root sense of
investigation or examination; Milton credits his colleague Peter Adamson
for the gloss \cite[p.~14]{Milton2011InductionBeforeHume}.

\paragraph{Al-F\={a}r\={a}b\={\i}.}
Al-F\={a}r\={a}b\={\i} (c.~870--c.~950) gave the first systematic Arabic
account. Milton reports Lameer's reading: for Aristotle induction advances
from related particular cases to the universal, whereas al-F\={a}r\={a}b\={\i}
explains it as an examination of the particulars. Lameer takes this to
follow from the Arabic rendering of \emph{epagoge} as \emph{istiqra'},
understood as a scrutiny of particulars. A consequence is that an
induction has full probative force only when complete
\cite[p.~14]{Milton2011InductionBeforeHume}.

\begin{hbdefinition}{def:p15-w02b-p15-026-tajriba}{Methodic experience (\emph{tajriba})}
Al-F\={a}r\={a}b\={\i} distinguishes induction from methodic experience
(\emph{tajriba}, corresponding to Greek \emph{empeiria})
\cite[pp.~14--15]{Milton2011InductionBeforeHume}.
\begin{itemize}
\item \emph{Methodic experience} examines the particular instances of a
  universal premise to see whether the universal is predicable of each,
  following this through all or most of them until necessary certainty is
  reached; the predication then applies to the whole species.
\item \emph{Induction} resembles it, but does not produce necessary
  certainty by universal predication, whereas methodic experience does.
\end{itemize}
Induction is accordingly the inferior: it neither uncovers necessary
truths nor yields certain knowledge.
\end{hbdefinition}

\paragraph{Avicenna.}
Avicenna (Ibn S\={\i}n\={a}, 980--1037) draws the same distinction in the
\emph{Book of Demonstration} of \emph{The Cure} (I.~9, \S\S~12, 21).
Unless induction examines all relevant cases it yields only probable
belief \cite[p.~15]{Milton2011InductionBeforeHume}.

\begin{hbexample}{def:p15-w02b-p15-026-scammony}{Scammony and the purging of bile}
Methodic experience is like judging that scammony purges bile. Because the
effect is repeated many times, it ceases to count as chance, and the mind
judges that purging bile is characteristic of scammony, a concomitant
accident of it. To the objection that the cause is unknown, Avicenna
replies that much repetition shows the connection is not by chance, since
chance is not always or for the most part. Scammony has no choice, so the
effect is something it brings about by nature
\cite[p.~15]{Milton2011InductionBeforeHume}.

Milton glosses this in later vocabulary, marked as his translation of the
point: we know a priori that the world is full of law-like regularities,
and experience is needed only to show that a given apparent regularity is
one of them and not chance \cite[p.~15]{Milton2011InductionBeforeHume}.
\end{hbexample}

Even granting this, methodic experience does not give certainty, for two
reasons Avicenna himself states \cite[pp.~15--16]{Milton2011InductionBeforeHume}.
\begin{enumerate}
\item \emph{Over-wide generalization.} Some country may have scammony
  whose temperament or special property prevents it from purging. The
  judgment from experience therefore concerns the scammony familiar to us
  and perceived by us. It purges bile, by essence or by a nature in it,
  unless some obstacle opposes.
\item \emph{The reference-class puzzle.} If the only people ever
  perceived were Sudanese and all were black, would repetition not produce
  the conviction that all people are black? If it does not, why does one
  repetition produce belief and another not? If it does, methodic
  experience has produced a falsehood. Milton calls this effectively the
  modern white-swan problem.
\end{enumerate}
Avicenna's answer, which Milton calls rather opaque, distinguishes what
the experience is taken to be about. If the procreation in question is
procreation by black people, or by the people of one country, the
experience is valid. If it is taken to be procreation by any people, the
experience does not end with the observed instances, since people
without qualification are not limited to black people
\cite[p.~16]{Milton2011InductionBeforeHume}. The answer restricts the
conclusion to the class actually sampled. It gives no criterion for
recognizing in advance when the sampled class is narrower than the class
in the conclusion.

This work of Avicenna was never translated into Latin. Scammony
nevertheless became a stock scholastic example, probably through his
medical writings, which shaped medical education in the Latin West
(citing Weinberg) \cite[p.~16]{Milton2011InductionBeforeHume}.

\subsection{The Latin West}

Two definitions circulated. One is Boethius' progression from particulars
to universals. The other is rhetorical, descending from Cicero through
Victorinus and Alcuin, and does not mention universality. Alcuin defines
induction as speech that proves uncertain things through certain ones and
brings the unwilling to assent; Victorinus speaks of capturing the
interlocutor's assent by undoubted matters
\cite[p.~16, n.~9]{Milton2011InductionBeforeHume}. When the main body of
Aristotle was received in the twelfth and thirteenth centuries, first from
Arabic and then from Greek, university attention shifted to the logical
tradition. Discussion of induction then occurs in commentaries and
questions on the \emph{Analytics} and in general logic treatises and
textbooks, few of which dwell on it. Milton notes that little has been
written specifically on medieval induction and points to Weinberg (1965)
and Bos (1993) \cite[p.~17]{Milton2011InductionBeforeHume}.

\paragraph{Robert Grosseteste.}
Grosseteste's commentary on the \emph{Posterior Analytics}, written before
1230, follows B19 but takes its example from Avicenna
\cite[p.~17]{Milton2011InductionBeforeHume}.

\begin{hbexample}{def:p15-w02b-p15-026-grosseteste}{Grosseteste's scammony procedure}
\begin{enumerate}
\item The senses repeatedly observe that eating scammony is accompanied by
  discharge of red bile, without seeing that scammony draws the bile out.
  From constant observation of the two observable things, a third,
  unobservable thing begins to be formed (\emph{estimare}): that scammony is
  the cause withdrawing the bile.
\item Repeated perceptions are stored in memory. Reason then begins to ask
  whether things are as remembered sensation suggests, and this leads to
  experiment (\emph{ad experientiam}).
\item Scammony is given once every other cause able to purge red bile has
  been separated off and ruled out. When this has been done many times with sure
  exclusion of the other purgatives, reason forms the universal that all
  scammony of its nature withdraws red bile. This is how one passes from
  sensation to a universal experimental principle.
\end{enumerate}
\cite[pp.~17--18]{Milton2011InductionBeforeHume}
\end{hbexample}

Milton makes three observations \cite[pp.~17--18]{Milton2011InductionBeforeHume}.
First, step~1 looks more like proposing a causal hypothesis than
generalizing. Second, the conclusion is stronger than ``all scammony draws
out red bile'': it is that scammony does so of its nature (\emph{secundum
se}). Grosseteste would presumably reply that the inference is safe only
if the power really belongs to the nature of scammony. Third, Grosseteste
departs from Aristotle in one important respect. In the \emph{Posterior
Analytics} many memories make one experience (\emph{empeiria}), which is
universal (100a5--6), and nothing like an experiment is required.
Grosseteste's procedure is interventionist: scammony is given in varied
situations from which the other known purgatives are excluded, and this
systematic variation justifies the universal conclusion.

\paragraph{William of Ockham.}
The most detailed account by a writer on logic is in Ockham's \emph{Summa
Logicae} III.~iii, chapters 31--36. Induction is defined, after Aristotle
and Boethius, as progression from singulars to a universal. The predicate
stays fixed and only the subject varies, as in ``this man runs, that man
runs, and so of the other singulars (\emph{et sic de singulis}); therefore
every man runs'', or the same with ``Socrates runs, Plato runs''. All
Ockham's examples predicate of individuals (Socrates, this man, that white
thing), never of species. Milton connects this with Ockham's nominalism:
only individuals exist, and universals are signs representing them
\cite[p.~18]{Milton2011InductionBeforeHume}.

\begin{hbdefinition}{def:p15-w02b-p15-026-ockham-rules}{Ockham's three rules}
For propositions \emph{de praesenti et de inesse}, Ockham gives three
rules \cite[p.~18]{Milton2011InductionBeforeHume}:
\begin{enumerate}
\item every true universal proposition has some true singular;
\item if all the singulars of a universal proposition are true, the
  universal is true;
\item if a negative universal proposition is false, at least one of its
  singulars is false.
\end{enumerate}
\end{hbdefinition}

Milton comments that rule~1 marks a basic difference from post-Fregean
logic, where ``every $A$ is $B$'' does not imply ``some $A$ is $B$''. Rule~2
might look obvious, but the later chapters show propositions for which
Ockham did not think it holds
\cite[pp.~18--19]{Milton2011InductionBeforeHume}. For modal propositions
in the divided sense (\emph{in sensu divisionis}), where the mode applies
to part of the proposition, the same rules hold. ``Socrates is contingently
an animal, Plato is contingently an animal, and so of the other singulars;
therefore every man is contingently an animal'' is accepted. In the
composite sense, where the mode applies to the whole proposition, they do
not \cite[p.~19, n.~11]{Milton2011InductionBeforeHume}:
\begin{itemize}
\item ``all the singulars are necessary, therefore the universal is
  necessary'' is not generally true, nor is its converse from universal
  to singulars;
\item ``all the singulars are possible, therefore the universal is
  possible'' does not hold (\emph{non valet}). Each of ``this contingent
  proposition is true'', ``that contingent proposition is true'', and so on
  may be possible without ``every contingent proposition is true'' being
  possible;
\item ``the singulars are contingent, therefore the universal is
  contingent'' does not hold (ch.~36).
\end{itemize}

\begin{hblemma}{thm:p15-w02b-p15-026-ockham}{proved=annex:p15-w02b-p15-026-ockham}{Ockham's rules, reconstructed}
Let $D$, $S$ and $P$ be as in \cref{thm:p15-w02b-p15-026-enumeration}.
The singulars of a universal with subject $S$ are the propositions
``$s$ is $P$'' (affirmative) or ``$s$ is not $P$'' (negative) for
$s\in S$. Read ``every $S$ is $P$'' with existential import, as true iff
$S\neq\varnothing$ and $S\subseteq P$, and read ``no $S$ is $P$'' as true
iff $S\cap P=\varnothing$.
\begin{enumerate}
\item Rule~1 holds. Under the modern reading of ``every $S$ is $P$'' as
  $S\subseteq P$ alone, it fails exactly when $S=\varnothing$.
\item Rule~2 holds whenever $S\neq\varnothing$. If $S=\varnothing$, every
  singular is vacuously true while the universal affirmative is false
  under the import reading.
\item Rule~3 holds.
\item Consider the composite possibility claim under the following
  hypotheses. A proposition is contingent iff it can be true and can be
  false. The contradictory of a contingent proposition is contingent. No
  proposition can be true together with its contradictory. Which
  propositions are contingent does not vary between the alternatives
  considered. At least one proposition is contingent. Then every singular
  ``this contingent proposition is true'' is possible, but ``every
  contingent proposition is true'' is not.
\end{enumerate}
\end{hblemma}

The lemma confirms Milton's verdict on these chapters. Their subject is
the deductive relation between universally quantified propositions and
their singulars. When Ockham says the inference from contingent singulars
to a contingent universal fails, he means \emph{all} the singulars.
Generalization from a finite sample is not raised
\cite[p.~19]{Milton2011InductionBeforeHume}. Part~(2) of the lemma is a
feature of the reconstruction. Milton reports nothing about how Ockham
treats subjects with no supposita. Existential import is treated at length
in \cref{sec:p02-w02-p02-001-historical-contrast}.

\paragraph{Jean Buridan.}
Buridan (c.~1300--1358), a nominalist who did comment on the
\emph{Analytics}, held that an induction is not a good consequence
(\emph{bona consequentia}) unless all the singulars are enumerated, which
is impossible because they are infinitely many. Yet we do generalize from
finite samples. Someone who has often seen rhubarb purge bile in varied
circumstances, remembers it and has never found a counter-example assents
to the universal. The assent comes from the intellect's natural inclination
to truth, not from a necessary consequence. The intellect supplies the
clause ``and so of the other singulars'', because it has never witnessed a
counter-example and sees no reason or dissimilarity why there should be one
\cite[pp.~19--20]{Milton2011InductionBeforeHume}.

Milton marks the difference from Hume that a modern reader might miss.
Buridan's inclinations are inclinations \emph{to the truth}, not habits
grounded on association of ideas. Behind them stands the unquestioned
assumption, notoriously absent in Hume, that God has given us faculties
which, properly used, lead to truth rather than error
\cite[p.~20]{Milton2011InductionBeforeHume}. In terms of
\cref{thm:p15-w02b-p15-026-enumeration}, the intellect's clause supplies
the missing exhaustiveness premise. Part~(3) of the lemma shows that this
supply is not a deduction from the observed cases.

\section{The Renaissance}
\label{sec:p15-w02b-p15-026-renaissance}

\subsection{The revival of rhetoric}

Milton allows that it simplifies matters, but says the Renaissance raised
rhetoric and lowered formal logic, unevenly and incompletely. Aristotle
kept respect, but the late medieval refinements, with what humanists
regarded as barbarous grammar and terminology, did not
\cite[p.~20]{Milton2011InductionBeforeHume}. Hostility goes back at least
to Petrarch. The first sustained attack was by Lorenzo Valla (1407--1457).
The earliest textbook in the new style was Rudolph Agricola's (1444--1485)
\emph{De Inventione Dialectica}, published in 1515 and often reprinted,
which Monfasani calls the first humanist logic to become a best seller.
Juan Luis Vives (1492--1540), who had studied logic in Paris and disliked
it, published \emph{In Pseudodialecticos} in 1520
\cite[p.~20]{Milton2011InductionBeforeHume}.

\begin{itemize}
\item \emph{Valla} opposed formal logic itself, not just its late medieval
  subtleties. Rules of reasoning, like rules of good writing, were to be
  drawn from the usage of the best Latin writers (\emph{ad consuetudinem
  eruditorum atque elegantium}), so logic became a minor part of rhetoric.
  He explicitly rejected the Boethian definition of induction as progression
  to universals. What mattered to him was the rhetorical argument from
  particulars to particulars \cite[pp.~20--21]{Milton2011InductionBeforeHume}.
\item \emph{Agricola} preferred \emph{enumeratio} to \emph{inductio}, both
  being Ciceronian, because Cicero had called it argument from enumerating
  all the parts. Some of his examples are traditional inductions, but
  others are not. One runs: ``the wall is mine, the foundation is mine, the
  roof is mine, the rest of the parts are mine. Therefore the house is
  mine.''
  \cite[p.~21]{Milton2011InductionBeforeHume}.
\item Arguments from parts to a whole seem to have become an accepted
  variety of rhetorical induction. Early in the eighteenth century Vico's
  \emph{Institutiones Oratoriae} distinguished \emph{inductio partium}
  from \emph{inductio similium}. The former in turn enumerates either all
  the species composing a genus or all the parts of a totality, such as
  the limbs and organs of a body \cite[p.~21]{Milton2011InductionBeforeHume}.
\end{itemize}

\subsection{Zabarella}

Jacopo Zabarella (1533--1589) taught at Padua, then the leading Italian
university and arguably the leading European one, where logic still
flourished. Schmitt calls him the most acute and influential of the Italian
Renaissance Aristotelians in methodological matters
\cite[p.~21]{Milton2011InductionBeforeHume}.

\begin{hbdefinition}{def:p15-w02b-p15-026-zabarella}{Dialectical and demonstrative induction (\emph{De Regressu}, ch.~4)}
Zabarella distinguishes two kinds of induction
\cite[pp.~21--22]{Milton2011InductionBeforeHume}.
\begin{enumerate}
\item \emph{Dialectical induction} is used in mutable and contingent
  matter (\emph{in materia mutabili et contingente}). It has no strength
  (\emph{nil roboris habet}) unless all the particulars without exception
  are considered (col.~485d).
\item \emph{Demonstrative induction} is used in necessary matter, among
  things with an essential connection. Not all the particulars are
  considered: having inspected some, the mind at once notices the essential
  connection, leaves the remaining singulars aside and infers
  (\emph{colligit}) the universal, knowing that the rest must be the same
  (col.~485d--e). A similar account is in \emph{De Methodis} III.~14
  (col.~255f).
\end{enumerate}
\end{hbdefinition}

For Aristotelians of Zabarella's kind, demonstrative induction was needed
because only it could give certain knowledge of the universal premises of
demonstrative syllogisms. Such truths come to be known not by an
impossible complete survey but by inspecting enough particulars for the
universal to form in the soul. What forms there is a real universal and not only a
concept: a de-individuated form, separated from matter
\cite[p.~22]{Milton2011InductionBeforeHume}. Milton represents the
situation with a diagram, reproduced here in tabular form.
\begin{center}
\begin{tabular}{ll}
(a) singular propositions & (b) universal propositions\\
(c) real individuals & (d) real universals
\end{tabular}
\end{center}
Logically, induction runs from (a) to (b); everyone in the Aristotelian,
as opposed to the rhetorical, tradition agreed on that. For medieval and
later realists, Zabarella included, the step from (a) to (b) is mirrored
by the relation of (c) to (d). What makes a universal proposition true is
what later philosophers might call a universal fact. That explains how a
universal can be known when not all particulars have been surveyed,
sometimes only a few (\emph{aliqua pauca}, col.~255f). Once the
intellect grasps the universal, no further examination of particulars is
needed, and in demonstrative induction certainty is attainable
\cite[p.~22]{Milton2011InductionBeforeHume}.

\begin{proofboundary}{def:p15-w02b-p15-026-realism}{Presuppositions of demonstrative induction}
Zabarella's account presupposes a realist theory of universals. Aristotle
apparently held one, in a form still disputed, and his medieval and later
successors developed several mutually incompatible ones (Milton cites his
own 1987 paper). The account was unavailable to nominalists such as Ockham,
for whom class (d) is empty. Nominalism remained a minority view among
university Aristotelians, though it became much more popular in the
seventeenth century \cite[p.~22]{Milton2011InductionBeforeHume}.
In terms of \cref{thm:p15-w02b-p15-026-enumeration}, demonstrative
induction does not claim that the inspected cases deductively settle
the universal. It claims a further, non-enumerative act by which the
intellect grasps the essential connection. Part~(3) of the lemma neither
supports nor refutes that claim. The claim is philosophical, not a result
about the reconstruction.
\end{proofboundary}

\section{Seventeenth and early eighteenth centuries}
\label{sec:p15-w02b-p15-026-modern}

Many of the most original philosophers of the period said little or
nothing about induction. The word does not occur in Spinoza's
\emph{Ethics} or Locke's \emph{Essay}, and appears once, in passing, in
Berkeley's \emph{Principles} \S~50. Spinoza does use it in chapter~11 of the
\emph{Tractatus Theologico-Politicus}. Locke uses it once, in \emph{The
Conduct of the Understanding}, where observations become reliable rules
when justified by a sufficient and wary induction of particulars. Milton
suggests a partial explanation of Locke's near-silence. Locke doubted that
natural philosophy would become a science, and as a physician he reasoned
about particular cases with general rules as fallible guides
\cite[p.~23, nn.~13--14]{Milton2011InductionBeforeHume}.

\subsection{Bacon}
\label{sec:p15-w02b-p15-026-bacon}

Francis Bacon was the first to invert the traditional order and put
induction before deduction. Milton nevertheless warns against calling him
the founder of inductive logic. Bacon was not a logician by temperament or
doctrine and is not usefully seen as a remote precursor of Carnap. His
induction belongs to a massive, unfinished programme for a new kind of
scientific knowledge (citing Malherbe and Gaukroger). Bacon and Descartes
were both dissatisfied with university natural philosophy, but Descartes
objected to its uncertainty while Bacon despised its uselessness for
controlling nature \cite[p.~23]{Milton2011InductionBeforeHume}.

\paragraph{The critique.}
In the \emph{Advancement of Learning} (1605) Bacon calls the logicians'
induction, by which the principles of sciences are supposedly found,
utterly vicious and incompetent. To conclude on an enumeration of
particulars without a contradictory instance is, he says, a conjecture and
no conclusion, for no one can be sure that unseen contrary instances do not
exist. What he would put in its place is not explained there; he promises
an account of his method, the Interpretation of Nature
\cite[pp.~23--24]{Milton2011InductionBeforeHume}. The account came in the
\emph{Novum Organum} (1620), the second, but first published, part of the
unfinished six-part \emph{Instauratio Magna}. Its title challenged
Aristotle's \emph{Organon} and was meant to replace it, not supplement it
\cite[p.~24]{Milton2011InductionBeforeHume}.

Milton identifies two strands in Bacon's case against school logic
\cite[pp.~24--25]{Milton2011InductionBeforeHume}.
\begin{enumerate}
\item \emph{Talk rather than action.} Bacon's science differs from
  ordinary logic in end, order of demonstration and starting point. It aims
  to invent arts rather than arguments and principles rather than things
  agreeing with principles, and to command nature in action rather than
  overcome an opponent in argument. Training in disputation prizes
  subtlety, but nature's subtlety far exceeds that of sense and
  understanding (\emph{Novum Organum} I.~10).
\item \emph{Form hides matter.} A syllogism consists of propositions,
  propositions of words, and words are symbols of notions. If the notions
  are confused or abstracted too hastily, nothing built on them is firm, so
  the only hope is a true induction (I.~14). This is the first mention of
  induction in the \emph{Novum Organum} proper. Milton notes that no later
  passage explains how induction would rectify defective concepts. He also
  notes the contrast with Descartes: for Bacon, concepts and knowledge
  improve together by slow increments. Clear and distinct ideas are not
  fixed first.
\end{enumerate}
Bacon did not reject the syllogism entirely. He left it jurisdiction over
popular arts and matters of opinion, while using induction throughout in
dealing with the nature of things (also in a letter to Redemptus Baranzan,
30~June~1622). Milton calls him the opposite of an ordinary-language
philosopher. One of Bacon's basic complaints against Aristotle was that
he had merely ordered popular notions instead of replacing them. Bacon
disliked neologisms and kept the ancient terms while altering their uses
\cite[p.~25, n.~15]{Milton2011InductionBeforeHume}.

\begin{hbdefinition}{def:p15-w02b-p15-026-bacon-two}{Simple enumeration and Bacon's new induction (\emph{Novum Organum} I.~105)}
Bacon separates two forms of induction \cite[pp.~25--26]{Milton2011InductionBeforeHume}.
\begin{enumerate}
\item \emph{Induction by simple enumeration} Bacon calls childish. What it
  concludes is precarious and at the mercy of a single contradictory
  instance, and it usually rests on too few facts, taken only from those
  near at hand.
\item \emph{The new induction} is to be used for proving and discovering all
  axioms, lesser and middle as well as first. It analyses nature by proper
  rejections and exclusions, and concludes on the affirmative instances
  only after a sufficient number of negatives. Furnishing it will take more
  labour than has been spent on the syllogism.
\end{enumerate}
``Axioms'' here are important general principles, not axioms in the modern
mathematical sense; the term comes from Stoic logic (Milton's note~16).
\end{hbdefinition}

Milton stresses that Bacon did not intend to keep simple enumeration with
safeguards that make its conclusions more probable. He wanted it discarded
for something different. The new induction is \emph{eliminative}, relying on
rejections and exclusions, and no mass of merely confirming instances is
ever enough. Bacon did not think the method could be described in detail
before being used. The specimen inquiry in the \emph{Novum Organum} is
explicitly a First Vintage, a provisional interpretation
(\emph{interpretatio inchoata}, II.~20). The full account was reserved for
the \emph{Scientia Activa}, the final part of the \emph{Instauratio}, which
was never written or even begun \cite[p.~26]{Milton2011InductionBeforeHume}.

\begin{hbmethod}{def:p15-w02b-p15-026-exclusion}{Bacon's three stages and the first work of exclusion (\emph{Novum Organum} II.~10--16)}
The method has three stages \cite[p.~26]{Milton2011InductionBeforeHume}.
\begin{enumerate}
\item Compile a ``natural and experimental history'' of the nature being
  studied (the \emph{given nature}, for example heat).
\item Order the history in tables (II.~10--14; Bacon describes these first
  two stages at length).
\item Induction. Its first work, so far as the discovery of Forms is
  concerned, is to reject or exclude each of the several natures that
  \begin{enumerate}
  \item is absent in some instance where the given nature is present; or
  \item is present in some instance where the given nature is absent; or
  \item increases in some instance where the given nature decreases; or
  \item decreases where the given nature increases (II.~16).
  \end{enumerate}
\end{enumerate}
After due rejection and exclusion, Bacon says, there will remain an
affirmative Form, solid, true and well defined, with light opinions gone
to smoke. Milton compares the image to gold left at the bottom of an
alchemist's crucible and quotes Bacon's own warning that this is quickly
said but the way to it is winding and intricate
\cite[pp.~26--27]{Milton2011InductionBeforeHume}.
\end{hbmethod}

Bacon's Forms are notoriously obscure and are not the substantial forms
of the Aristotelians. Whatever they are ontologically, they are the causes
of the phenomenal natures (citing P\'erez-Ramos). The form of heat is
present in all hot bodies, absent from all cold ones, and varies in
intensity with the body's degree of heat
\cite[p.~26]{Milton2011InductionBeforeHume}.

\begin{hbdefinition}{def:p15-w02b-p15-026-limited-variety}{The Limited Variety presupposition}
Bacon's confidence that eliminative induction yields certainty rested on
several presuppositions. Milton names the most important as what Keynes
later called a Principle of Limited Variety. The world's apparently endless
variety arises from combinations of a finite, indeed fairly small, number
of simple natures, which compose an alphabet of nature (Bacon's
\emph{Abecedarium Naturae}) which cannot be guessed or found by
speculation but begins to appear once correct inductive procedures are
used \cite[p.~27, n.~17]{Milton2011InductionBeforeHume}.
\end{hbdefinition}

Bacon offered no a priori justification of this presupposition. Milton
sees no reason to think Bacon would have regarded one as possible or
needed. Validation would be retrospective: the method would be vindicated
by giving those who used it correctly power over nature
\cite[p.~27]{Milton2011InductionBeforeHume}. The following reconstruction
states, in modern terms, what the exclusion clauses deliver when
presuppositions of this kind are \emph{assumed}. It is not a justification
Bacon gave, and it does not establish the presuppositions.

\begin{hblemma}{thm:p15-w02b-p15-026-exclusion}{proved=annex:p15-w02b-p15-026-exclusion}{Exclusion under limited variety}
Let $I$ be a finite set of recorded instances, and let $\mathcal N$ be a set
of candidate natures. For the given nature $g$ and each candidate
$n\in\mathcal N$ let $\delta_g,\delta_n\colon I\to[0,\infty)$ give degrees,
with presence meaning a positive degree. Exclude $n$ if for some
$i,j\in I$:
\begin{itemize}
\item[(E1)] $\delta_g(i)>0$ and $\delta_n(i)=0$; or
\item[(E2)] $\delta_g(i)=0$ and $\delta_n(i)>0$; or
\item[(E3)] $\delta_g(j)<\delta_g(i)$ and $\delta_n(j)>\delta_n(i)$.
\end{itemize}
(Clause E3 covers both of Bacon's covariation clauses: exchanging $i$ and
$j$ turns one into the other.) Assume the \emph{form hypothesis}: some
$f\in\mathcal N$ satisfies $\delta_f(i)>0\iff\delta_g(i)>0$ for every
$i\in I$, and $\delta_g(j)<\delta_g(i)\Rightarrow\delta_f(j)\leq\delta_f(i)$
for all $i,j\in I$. Then:
\begin{enumerate}
\item $f$ is not excluded;
\item if exactly one candidate survives exclusion, it is $f$;
\item if two candidates satisfy the form hypothesis's conditions on $I$,
  both survive, and the instances do not identify the form;
\item if no candidate satisfies the form hypothesis, a candidate may be the
  unique survivor without being a form of $g$.
\end{enumerate}
\end{hblemma}

Parts~(1)--(2) show what the presuppositions buy. If the form is among
finitely listed candidates and the records are accurate, elimination is
conclusive whenever it leaves one survivor. Parts~(3)--(4) show where the
weight lies. Without enough discriminating instances the procedure stops
short, and without the hypothesis that the cause is in the list the
survivor proves nothing. The finiteness of $\mathcal N$ is what makes a
unique survivor attainable from finitely many instances. That finiteness is
Milton's Limited Variety, and he records that Bacon gave it no a priori
defence. The covariation relation used here is the one later stated by
Mill (\cref{def:p15-wm-variation}).

\subsection{Descartes}

Induction plays no significant role in Descartes's mature philosophy, but
the early, unfinished \emph{Regulae ad Directionem Ingenii}
(c.~1619--c.~1628) contains remarks on it. Whether Descartes had then read
Bacon is unknown. He later knew Bacon's thought, but the \emph{Regulae}
account appears to owe nothing to the \emph{Novum Organum}. In the
\emph{Regulae}, intuition, a direct apprehension unmediated by other
operations, gives the most certain knowledge. Deduction is needed because
some chains are too long to grasp in one act: each link can be intuited,
but not all at once \cite[p.~27]{Milton2011InductionBeforeHume}.

\begin{hbdefinition}{def:p15-w02b-p15-026-sufficient}{Sufficient enumeration (Rule~VII)}
Rule~VII requires that, for knowledge to be complete, everything bearing
on the undertaking be surveyed in a continuous, uninterrupted movement of
thought and included in a sufficient and well-ordered enumeration
(\emph{sufficienti et ordinata enumeratione})
\cite[p.~27]{Milton2011InductionBeforeHume}.
\begin{enumerate}
\item \emph{Name.} For Descartes \emph{enumeratio} and \emph{inductio}
  seem to be two names for one thing; he writes ``enumeratio, sive
  inductio''. Milton's note~18 counts three occurrences of
  \emph{inductio} in Rule~VII and one in Rule~XI. A further occurrence in
  Rule~III of the 1701 edition may be an error for \emph{deductio}
  \cite[p.~27]{Milton2011InductionBeforeHume}.
\item \emph{Function.} Enumeration completes knowledge. It is a survey of
  everything relevant to the problem, careful enough that we may conclude
  with manifest certainty that nothing has been overlooked. Even when the
  object eludes us, an enumeration shows with certainty that it cannot be
  found by any method known to us \cite[p.~28]{Milton2011InductionBeforeHume}.
\item \emph{Sufficiency.} A sufficient enumeration, or induction, is one
  that makes a conclusion more certain than any other kind of proof except
  simple intuition allows. When knowledge is not reducible to intuition and
  syllogistic fetters are cast off, it is the one remaining path
  \cite[p.~28]{Milton2011InductionBeforeHume}.
\end{enumerate}
\end{hbdefinition}

Milton finds Descartes's explanation of sufficiency regrettably short of
his own ideal of clarity (citing Beck). Sufficiency is not completeness
\cite[pp.~28--29]{Milton2011InductionBeforeHume}.

\begin{hbexample}{def:p15-w02b-p15-026-descartes-cases}{Descartes's three enumerations}
\begin{enumerate}
\item \emph{Corporeal kinds and the rational soul.} To determine how many
  kinds of corporeal thing there are, every kind must be distinguished and
  the enumeration must be complete. To show that the rational soul is not
  corporeal, completeness is unnecessary. It suffices to group all bodies
  into several classes and show the soul belongs to none. Milton reads
  this as follows: a negative thesis about what the soul is not needs only
  a division into a few broad classes \cite[pp.~28--29]{Milton2011InductionBeforeHume}.
\item \emph{Paths to truth} (Rule~VIII). Whoever investigates all kinds of
  knowledge begins with the pure intellect and then enumerates the other
  instruments of knowledge, of which there are only two, imagination and
  sense perception. The paths to truth open to us are few enough to be
  found at once by a sufficient enumeration
  \cite[p.~29]{Milton2011InductionBeforeHume}.
\item \emph{The isoperimetric example.} To show by enumeration that a
  circle has greater area than any other figure of equal perimeter, one
  need not examine every figure. Proving it for some particular figures
  entitles one to conclude by induction for all the rest. Milton finds this
  puzzling. It is not a proof by an exhaustive finite case division, like a
  proof for triangles split into acute, obtuse and right-angled. There are
  infinitely many polygons with the circle's perimeter and smaller area.
  How the argument is meant to go is unclear, but it is not by exhaustive
  cases. His note~19 adds that the result is not elementary: a non-rigorous
  proof by Zenodorus (second century BC?) is preserved in Book~V of Pappus'
  \emph{Collection} \cite[p.~29]{Milton2011InductionBeforeHume}.
\end{enumerate}
\end{hbexample}

Case~1 is part~(2) of \cref{thm:p15-w02b-p15-026-enumeration}. If the
classes $K_1,\dots,K_m$ jointly exhaust body and the soul is in none, the
soul is not a body, however coarse the classes. The deductive force comes
from the covering premise, which is far easier to secure than a complete
list of kinds. Case~3 has no such covering premise. By part~(3) of the lemma,
the checked figures alone do not settle the universal, which agrees with
Milton's verdict that the argument is not an exhaustive case analysis. The
lemma does not prove the isoperimetric theorem, and the theorem is not
stated in this chapter.

The \emph{Regulae} were first printed in 1701, though manuscript copies
were read by Leibniz and possibly by Locke. If Locke saw one, it would have
been during his Dutch exile of 1683--1689, and his papers do not mention
it. Book~IV of Locke's \emph{Essay} has close parallels with the
\emph{Regulae} but never mentions induction
\cite[p.~29, n.~20]{Milton2011InductionBeforeHume}.

\subsection{Gassendi}

Pierre Gassendi's \emph{Institutio Logica}, published in 1658 in his
posthumous \emph{Opera Omnia}, follows the \emph{Prior Analytics} in
treating induction as a kind of syllogism
\cite[pp.~29--30]{Milton2011InductionBeforeHume}.

\begin{hbexample}{def:p15-w02b-p15-026-gassendi}{Gassendi's concealed premise}
Every walking animal lives, every flying animal lives, and also every
swimming, creeping and plant-like animal; therefore every animal lives.
The induction has a concealed premise: every animal is walking, flying,
swimming, creeping or plant-like. Gassendi says that without it the consequence would
have no force (\emph{consequutionis vis nulla foret}): were some further
kind of animal to exist, the conclusion could come out false. A legitimate
(\emph{legitima}) induction must rest on an enumeration of all relevant
species or parts, and Gassendi remarks that such an enumeration is usually
difficult if not impossible \cite[p.~30]{Milton2011InductionBeforeHume}.
Bernier's French \emph{Abr\'eg\'e} (1684) repeats the account with minor
changes \cite[p.~30]{Milton2011InductionBeforeHume}.
\end{hbexample}

In the terms of \cref{thm:p15-w02b-p15-026-enumeration}, Gassendi's
concealed premise is the covering condition of part~(2), with five
classes, and his remark about a further kind of animal is part~(3) applied
to an unlisted class. The same structure is Aristotle's B23 inference read
over species (\cref{thm:p15-w02b-p15-026-conversion}).

\subsection{Arnauld and Nicole}

The work most influenced by the then unpublished \emph{Regulae} was
Arnauld and Nicole's \emph{La Logique ou l'art de penser} (1662). It
introduces induction neutrally. Induction occurs when examining several
particular things leads us to a general truth. Having found several seas
salty and several rivers fresh, we infer that sea water in general is salty
and river water fresh. Induction is called the beginning of all knowledge,
because singulars are presented before universals
\cite[p.~30]{Milton2011InductionBeforeHume}.

Milton finds the Aristotelian resemblance superficial. One might never have
thought about triangles without seeing one, yet the general and certain
conclusion about all triangles rests on considering the idea of a triangle
in one's mind, not on examining all triangles. The same holds for very
general axioms, such as Euclid's last Common Notion that the whole is
greater than its part. Unnamed philosophers, presumably Gassendi and his
followers, said we know this only from childhood observation that a man is
larger than his head, a house than a room, a forest than a tree. Arnauld
and Nicole reply that we would then be sure only of its probability, since
induction is a certain means of knowledge only when complete
\cite[pp.~30--31]{Milton2011InductionBeforeHume}.

Several of their examples of over-confident induction come from recent
physical science \cite[p.~31]{Milton2011InductionBeforeHume}.
\begin{itemize}
\item Natural philosophers long held that a piston cannot be drawn from a
  perfectly sealed syringe and that a suction pump can lift water from any
  depth. They believed both rested on a very certain induction from an
  infinity of experiments.
\item It was assumed that water in a curved vessel such as a U-tube stands
  at equal heights in both arms even when one arm is wider. Experiment
  showed otherwise when one arm is very narrow and capillary attraction
  matters.
\end{itemize}
Milton's assessment is guarded. It would probably go too far to call new
scientific discoveries the main force behind inductive scepticism, but they
seem to have helped weaken trust in the long-accepted experiential data
underlying Aristotelian science (citing Dear)
\cite[p.~31]{Milton2011InductionBeforeHume}. These are historical examples
reported in a logic textbook, not measurements Milton reanalyses.

\subsection{Hobbes and Wallis}

Despite Hobbes's firm nominalism, induction plays little part in his
philosophy. In the mid-1650s, in acrimonious exchanges with the
mathematician John Wallis (citing Jesseph), he attacked Wallis's use of
induction. In the \emph{Six Lessons} (1656) Hobbes ridicules the idea that an
induction which does not number every particular could license a universal
conclusion and pass for geometrical demonstration. Milton infers
that Hobbes recognized just two kinds of induction: one founded on a
complete enumeration, which can give certainty, and one founded on a
partial enumeration, which cannot \cite[p.~31]{Milton2011InductionBeforeHume}.

Wallis's remarks were in the \emph{Arithmetica Infinitorum} (1656). His
fullest defence is in the much later \emph{Treatise of Algebra} (1685),
answering a far more capable critic, Fermat (citing Stedall). Wallis
holds that induction is a very good method of investigation: it often
leads easily to a general rule, or at least prepares for one. When the
inquiry reveals a regular, orderly progression proceeding by one general
process, with no ground to suspect failure, the result needs no further
demonstration, though it may admit one
\cite[pp.~31--32]{Milton2011InductionBeforeHume}.

\begin{hbexample}{def:p15-w02b-p15-026-wallis-binomial}{Wallis's binomial table}
Wallis's example is the expansion of the binomial $(a+e)^n$:
\begin{align*}
(a+e)^2&=a^2+2ae+e^2,\\
(a+e)^3&=a^3+3a^2e+3ae^2+e^3,\\
(a+e)^4&=a^4+4a^3e+6a^2e^2+4ae^3+e^4,
\end{align*}
and so on. The coefficients can be found by repeated multiplication, but
more easily from the arrangement now called Pascal's triangle, in which each
entry is the sum of the two diagonally above it. Milton prints the rows
$1$; $1\ 1$; $1\ 2\ 1$; $1\ 3\ 3\ 1$; $1\ 4\ 6\ 4\ 1$;
$1\ 5\ 10\ 10\ 5\ 1$ \cite[p.~32]{Milton2011InductionBeforeHume}.

That the procedure works for every power is said to be established by
induction. Wallis notes that most mathematicians, after continuing such an
induction a few steps and seeing no reason to doubt that it continues,
conclude universally, and that such induction has been thought a
conclusive argument by those not inclined to be captious. He gives the
reason: the nature of Number contains a sufficient ground for the sequel
\cite[p.~32]{Milton2011InductionBeforeHume}.
\end{hbexample}

Milton draws three boundaries here \cite[p.~32]{Milton2011InductionBeforeHume}.
Wallis never suggests that results reached this way are merely probable.
Wallis is \emph{not} using what is now called mathematical induction
(citing Cajori). His argument is closer to the demonstrative induction of
later Aristotelians such as Zabarella, in which a few cases reveal the
underlying regularity. The following lemma states the modern fact about the
table. Its proof is a modern argument and should not be read back into
Wallis.

\begin{hblemma}{thm:p15-w02b-p15-026-pascal}{proved=annex:p15-w02b-p15-026-pascal}{The table rule for every power}
For integers $n\geq 0$ and $0\leq k\leq n$, let $c(n,k)$ be the coefficient
of $a^{n-k}e^k$ in the expansion of $(a+e)^n$ in commuting indeterminates
$a,e$, and put $c(n,k)=0$ for $k<0$ or $k>n$.
\begin{enumerate}
\item For every $n\geq 0$ and every integer $k$,
  $c(n+1,k)=c(n,k-1)+c(n,k)$.
\item Let $t(n,k)$ be the table built from $t(0,0)=1$ (zero outside
  $0\leq k\leq n$) by the rule $t(n+1,k)=t(n,k-1)+t(n,k)$. Then
  $t(n,k)=c(n,k)$ for all $n\geq 0$ and all $k$.
\end{enumerate}
\end{hblemma}

Part~(1) is a single identity holding for each $n$, obtained by
multiplying once more by $a+e$. Part~(2) is the step that needs
mathematical induction on $n$. Neither Wallis nor Milton states
anything like part~(1), and Milton does not say what Wallis's ``ground in
the nature of Number'' consists in. The lemma only shows that the modern
justification has this two-part shape.

Wallis's \emph{Institutio Logicae} (1687), Book~III, chapter~15, is more
traditional and largely concerned with reducing perfect inductions to
syllogistic figures. The conclusion of an imperfect induction is only
conjectural or probable, because a single negative instance can overturn
it. No degree of probability is estimated. Wallis insists that the
weakness (\emph{imbecillitas}) of imperfect induction lies in its matter,
not its form, which is that of a syllogism
\cite[p.~33]{Milton2011InductionBeforeHume}.

\begin{hbexample}{def:p15-w02b-p15-026-teeth}{Wallis's horned animals}
Someone argues, in Milton's translation of Wallis: ``Teeth in the upper
jaw are absent in all horned animals; because it is thus in the Ox, the
Sheep, the Goat, nor is it otherwise (as far as we know) in the others;
Therefore (at least as far as we know) in all.'' The conclusion is only probable
(\emph{verisimilis}), not through a defect of syllogistic form but through
the uncertainty of the matter, that is, of the premises
\cite[p.~33]{Milton2011InductionBeforeHume}.

Milton's reading: on this account an inductive argument is not a fallible
inference from secure premises such as ``the ox has no upper teeth''. It is
a deductive inference from premises such as ``the ox, sheep and goat lack
upper teeth, and it is not otherwise in the other horned animals''. Those
premises are all uncertain and are accepted provisionally because no
counter-example is known \cite[p.~33]{Milton2011InductionBeforeHume}.
\end{hbexample}

Wallis's reading places the uncertainty in the covering premise of
part~(2) of \cref{thm:p15-w02b-p15-026-enumeration}. The inference is then
valid, and the doubt attaches to whether the listed cases together with
``the others'' exhaust the subject.

\subsection{Leibniz}
\label{sec:p15-w02b-p15-026-leibniz}

Induction was not central for Leibniz, but what he said is of interest
(citing Rescher and Westphal) \cite[p.~33]{Milton2011InductionBeforeHume}.

\paragraph{Truths of reason.}
The preface to the \emph{New Essays} names innate principles as a basic
point of disagreement with Locke, and with it the question whether every
truth rests on experience (induction and instances) or some have another
foundation. Leibniz takes the second option. The senses supply only
instances, that is, particular truths. No number of confirming instances
yields the universal necessity of a general truth, since what has happened
need not always recur in the same way. Arithmetic, geometry and the other
truths of reason do not rest on induction at all
\cite[p.~33]{Milton2011InductionBeforeHume}.

\paragraph{Empirical inference.}
Truths of fact must be found otherwise, at least by humans. In a paper on
the souls of men and beasts (around 1710) Leibniz distinguishes empirical
and rational inferences. Empirical inferences are shared with beasts. Having
sensed things repeatedly connected, we expect them connected again: dogs
beaten after doing something displeasing expect a beating if they repeat
it, and so do infants. \emph{Monadology} \S~28 adds that people act like
beasts so far as their perceptions follow memory alone, like empirical
physicians with practice but no theory. We are mere empirics in
three-quarters of our actions. Expecting daylight tomorrow because it has
always come is empirical; only the astronomer concludes it by reason. One
difference remains. So far as can be seen, beasts have no awareness that a
proposition is universal. Empirics are sometimes led by induction to
true universal propositions, but only by accident, not by the force of
consequence. Milton's note~21 records two further places where Leibniz
compares beasts with empirical physicians: \emph{Principles of Nature and
Grace} \S~5 and the preface to the \emph{New Essays}
\cite[pp.~33--34, n.~21]{Milton2011InductionBeforeHume}.

\begin{hbdefinition}{def:p15-w02b-p15-026-leibniz-grades}{Leibniz's grades of confirmation}
Milton infers that Leibniz recognizes three kinds of reasoning, in a broad
sense of the word: (1)~inference from particulars to particulars,
earlier called \emph{paradeigma} or example; (2)~inductive generalization
proper; (3)~deduction. He finds support in a note on the back of a draft
letter of May 1693 distinguishing three grades of confirmation
(\emph{firmitas}) \cite[p.~34]{Milton2011InductionBeforeHume}:
\begin{enumerate}
\item \emph{logical certainty}, found in propositions of eternal truth;
\item \emph{physical certainty}, which is only logical probability, found in
  propositions known by induction, such as that every man is a biped,
  though some are born with one foot or none;
\item \emph{physical probability}, as in the south wind bringing rain, which
  is usually true but often enough false.
\end{enumerate}
Physical certainty is identified with moral certainty in \emph{New
Essays} IV.~vi.~13. Milton concludes that inductive conclusions can, at
least sometimes, be morally certain \cite[p.~34]{Milton2011InductionBeforeHume}.
\end{hbdefinition}

\paragraph{The \emph{adminicula}.}
How such certainty is reached appears in one of Leibniz's earliest
treatments, the preface (1670) to his edition of Mario Nizzoli's
\emph{De Veris Principiis} (first published 1553). Nizzoli (1488--1567)
has been called a Ciceronian Ockhamist, and Leibniz takes his basic error
to be nominalism. If universals were only collections of individuals, there
could be no knowledge by demonstration, only by collecting individuals,
that is, by induction. Milton's note~22 points out that the Latin
\emph{collectionem singularium, seu inductionem} names one process twice,
not two processes. Knowledge would then be impossible and the sceptics would
win, because perfectly universal propositions can never be established that
way: one can never be sure all individuals have been considered, and must
stop at ``all the cases I have experienced are so''
\cite[p.~35]{Milton2011InductionBeforeHume}.

\begin{hbdefinition}{def:p15-w02b-p15-026-adminicula}{Leibniz's helping propositions (\emph{adminicula})}
We are confident that fire burns and that a hand placed in it will be
burned. That moral certainty does not rest on induction alone but on
induction aided by further universal propositions
\cite[p.~35]{Milton2011InductionBeforeHume}:
\begin{enumerate}
\item where the cause is the same or similar in every case, so is the
  effect;
\item a thing that is not sensed is not assumed to exist;
\item what is not assumed is ignored in practice until it has been
  proved.
\end{enumerate}
\end{hbdefinition}

Milton classifies the second and third as methodological principles,
similar to but not identical with Ockham's Razor. He calls the first a more
carefully worded form of Hume's principle ``that like causes always produce
like effects'', the sixth of Hume's rules for judging causes in the
\emph{Treatise} I.~iii.~15. Without these aids not even moral certainty is
possible. Their own knowledge therefore cannot rest on induction, since
inductively derived aids would need further aids, and so on to infinity.
Leibniz's word is Baconian: the \emph{Adminicula inductionis} were announced
in \emph{Novum Organum} II.~21 but never described. His thought, Milton
says, is much closer to Hume
\cite[p.~35, nn.~22--24]{Milton2011InductionBeforeHume}.

\begin{hbwarning}{def:p15-w02b-p15-026-zabarella-leibniz}{Two ways of supplementing induction}
Milton compares Leibniz's buttressed induction with Zabarella's
demonstrative induction \cite[p.~36]{Milton2011InductionBeforeHume}.
\begin{itemize}
\item In \emph{demonstrative induction}, certainty comes from the
  intellect's grasp of the universal nature that grounds the universal
  proposition.
\item In \emph{Leibniz}, help comes from principles of much greater
  generality, such as the Law of Continuity and ultimately the Principle of
  Sufficient Reason.
\end{itemize}
He quotes Foucher de Careil's verdict that, for Leibniz, induction needs
universal propositions independent of it. The law of continuity links
universal and particular and is the true basis of induction, without which
induction is sterile and with which it generates moral certainty. Milton's
summary is that the regularity of nature, its being law-governed, makes
properly conducted inductive inferences safe. The supplement is in each
case a philosophical thesis, not a consequence of the observed instances.
\end{hbwarning}

\section{Milton's conclusion}
\label{sec:p15-w02b-p15-026-conclusion}

The history is episodic and fragmentary. Its material is remarks
extracted from writers who were almost always concerned with other matters
and for whom induction was of minor importance. The one significant
exception is Bacon, and even his treatment of eliminative induction is
remarkably brief given its pivotal place in his programme
\cite[p.~36]{Milton2011InductionBeforeHume}. Milton adds that even now a
philosopher need not treat induction to be a candidate for greatness: no one
would have expected Nietzsche to focus on it. What is striking is not that
many chose other topics but that virtually everyone did
\cite[p.~36]{Milton2011InductionBeforeHume}.

Talk of a single, wide-ranging problem of induction is comparatively
new. Among the first and most influential uses of the phrase is the
peroration of Mill's \emph{System of Logic} III.~iii.~3. It asks why a
single instance sometimes suffices for a complete induction while myriads
of concurring instances, with no known exception, sometimes go a very
little way towards a universal proposition; whoever can answer that, Mill
says, has solved the problem of induction. Milton doubts that anyone has
since solved or dissolved it, though confident and sometimes absurd claims
continue. What does seem clear to him is that before the nineteenth century
induction was not seen as raising one general problem. Still less did anyone
regard failure to solve it as, in C.~D. Broad's often-quoted phrase,
``the scandal of Philosophy''
\cite[pp.~36--37]{Milton2011InductionBeforeHume}.

\section{Comparisons, limitations and relationships}
\label{sec:p15-w02b-p15-026-comparisons}

\Cref{tab:p15-w02b-p15-026-comparison} collects the positions above by what
induction is said to lead to and what certainty it is said to give. Each row
reports Milton's account at the pages cited in the corresponding section.

\begin{table}[ht]
\centering
\small
\begin{tabular}{p{0.2\linewidth}p{0.33\linewidth}p{0.37\linewidth}}
\hline
Writer & Direction and object & Certainty claimed or denied\\
\hline
Socrates, Plato & cases to a philosophical universal & graspable directly; not empirical\\
Aristotle & individuals to universals (concepts and/or propositions); example: part to part & B23 needs ``through all''; planets premise (2) fallible\\
Diogenes Laertius & true premises to a resembling conclusion & from opposites; from implication (rhetorical, dialectical)\\
Sextus & particulars to universal & some: insecure; all: impossible\\
Alexander & particulars to universal & persuasive, not necessary\\
Cicero, Quintilian & agreed cases to a similar case & judged for persuasiveness\\
al-F\={a}r\={a}b\={\i}, Avicenna & examination of particulars & induction certain only if complete; methodic experience certain\\
Grosseteste & observation, then experiment with exclusion & universal of the nature of the thing\\
Ockham & all singulars to universal & deductive rules; composite modal exceptions\\
Buridan & many singulars to universal & assent by natural inclination to truth\\
Valla, Agricola, Vico & particulars to particulars; parts to whole & rhetorical\\
Zabarella & few cases to universal (demonstrative) & certain in necessary matter, via real universals\\
Bacon & exclusion to an affirmative Form & certain given Limited Variety; simple enumeration childish\\
Descartes & sufficient, ordered enumeration & certain if sufficient; sufficiency $\neq$ completeness\\
Gassendi & species to genus & valid only with covering premise\\
Arnauld, Nicole & several cases to general truth & certain only if complete\\
Hobbes & complete or partial enumeration & certain only if complete\\
Wallis & few cases to rule (algebra); listed cases plus ``others'' (logic) & conclusive in algebra; probable in matter, valid in form\\
Leibniz & instances plus \emph{adminicula} & moral certainty; never universal necessity\\
\hline
\end{tabular}
\caption{Positions on induction before Hume, as Milton reports them.}
\label{tab:p15-w02b-p15-026-comparison}
\end{table}

\paragraph{Milton's comparisons across writers.}
Gideon is set against Mill's methods (p.~2). Aristotle's B19 is set against
Locke's abstract general ideas (p.~6). Sextus's crocodile is set against the
black swan and Lakatos's monster-barring (p.~10). Avicenna's Sudanese case
is set against the white-swan problem (p.~16). Grosseteste is set against
the non-experimental \emph{empeiria} of the \emph{Posterior Analytics}
(p.~18). Buridan's inclination to truth is set against Hume's habit (p.~20).
Bacon is set against Descartes and against Carnap (pp.~23, 25).
Descartes's \emph{Regulae} is set against Locke's \emph{Essay} IV (p.~29).
Wallis's algebraic induction is set against mathematical induction and
Zabarella (p.~32). Leibniz's \emph{adminicula} are set against Ockham's
Razor and Hume's rules, and his supplemented induction against Zabarella's
(pp.~35--36)
\cite[pp.~2--36]{Milton2011InductionBeforeHume}.

\paragraph{Limitations Milton states.}
Milton describes the history as patchy and episodic. A larger surviving
ancient literature would make it longer and fuller (p.~12). He does not
discuss Hellenistic sign inference in detail
(\cref{def:p15-w02b-p15-026-signs}). He does not analyse the enthymeme
(p.~7, note~4). He does not decide whether Hume was an inductive sceptic
(p.~1). He states his preference for Ross over Barnes as an inclination,
not a resolution (p.~6). He calls his explanation of Aristotle's lack of a
systematic account one possible explanation (p.~8)
\cite[pp.~1, 6--8, 12]{Milton2011InductionBeforeHume}.

\begin{proofboundary}{def:p15-w02b-p15-026-reconstruction-scope}{What the lemmas of this chapter establish}
\Cref{thm:p15-w02b-p15-026-conversion,thm:p15-w02b-p15-026-enumeration,thm:p15-w02b-p15-026-ockham,thm:p15-w02b-p15-026-exclusion,thm:p15-w02b-p15-026-pascal}
are mathematical statements about explicitly defined extensional, modal
or algebraic models. They make exact the separations Milton draws between
valid and invalid enumerative forms, between complete, covering and partial
surveys, between the deductive rules of Ockham's chapters and
generalization from samples, and between what elimination delivers with and
without its presuppositions. They do not show that any historical writer
used such a model. They assign no probability to an inductive conclusion.
They do not decide the interpretive disputes Milton records (Asmis and
Barnes, Ross and Barnes, Burnyeat's reading of example). They do not
establish the philosophical presuppositions on which the certainty claims
of Zabarella, Bacon and Leibniz depend.
\end{proofboundary}

\paragraph{Relationships.}
Hume's argument, whose novelty Milton contrasts with the commonplace
fallibility of enumeration, is reconstructed in
\cref{sec:p15-hume-and-the-problem-of-induction}. Mill's methods, the
background of Milton's Gideon comparison and the modern form of Bacon's
covariation clause, are in \cref{sec:p15-wm-methods}. The syllogistic
apparatus behind Aristotle's B23 and Gassendi's account is in
\cref{sec:p02-syllogistic-inference}. Existential import, the point of
Ockham's first rule, is in \cref{sec:p02-w02-p02-001-historical-contrast}.
Goodman's later problem of which generalizations are projectible is in
\cref{sec:p15-goodman-entrenchment-and-projectibility}.

\index{induction!history before Hume}
\index{epagoge}
\index{Aristotle!epagoge}
\index{paradeigma}
\index{methodic experience}
\index{Zabarella, Jacopo!demonstrative induction}
\index{Bacon, Francis!eliminative induction}
\index{limited variety}
\index{Descartes, Ren\'e!sufficient enumeration}
\index{Wallis, John!induction}
\index{Leibniz, Gottfried Wilhelm!adminicula}
